\chapter{Введение}
\label{ch:intro}

\section{Актуальность темы}

Городская среда предоставляет множество альтернативных пешеходных маршрутов между
двумя точками. Стандартные навигационные сервисы, как правило, минимизируют время
или расстояние, не учитывая субъективные предпочтения: прохождение через парки,
избегание шумных магистралей, выбор участков с лучшим качеством воздуха или
набережных. Для жителей мегаполисов прогулка --- не только способ перемещения, но
и досуг, поэтому задача \emph{персонализированного} планирования маршрута является
актуальной как с практической, так и с научной точки зрения.

Развитие открытых геоданных (OpenStreetMap), свободных маршрутизаторов (OSRM) и
современных веб-технологий позволяет создавать специализированные приложения без
лицензирования проприетарных картографических API. При этом математически задача
сводится к комбинации методов многокритериальной оценки объектов, построения
взвешенного графа и классического алгоритма поиска кратчайшего пути.

\section{Цель и задачи работы}

\textbf{Цель} --- разработать и реализовать веб-приложение Fun-Walk для
автоматизированного планирования пешеходных маршрутов с учётом шести
пользовательских критериев: парки, зелёные зоны, велодорожки, качество воздуха,
тихие места и набережные.

Для достижения цели поставлены следующие \textbf{задачи}:
\begin{enumerate}
  \item провести анализ предметной области и существующих подходов к маршрутизации;
  \item формализовать математическую модель оценки точек интереса (POI) и построения
        взвешенного графа маршрута;
  \item спроектировать архитектуру клиент-серверного приложения;
  \item реализовать алгоритм Dijkstra и интеграцию с OSRM на языке TypeScript;
  \item разработать пользовательский интерфейс с интерактивной картой;
  \item провести тестирование и оценить качество получаемых маршрутов.
\end{enumerate}

\section{Объект, предмет и методы исследования}

\textbf{Объект} --- процесс планирования пешеходных маршрутов в городской среде.

\textbf{Предмет} --- математические модели и программные алгоритмы, обеспечивающие
выбор последовательности точек интереса и построение проходимой геометрии маршрута.

\textbf{Методы:} системный анализ; линейная алгебра; геодезия на сфере; теория
графов; алгоритм Dijkstra; объектно-ориентированное программирование; REST API;
функциональное тестирование.

\section{Практическая значимость}

Результат работы представляет собой работоспособный прототип, демонстрирующий
применение математических методов в прикладной разработке. Код проекта может
служить основой для расширения: подключение реальных API мониторинга воздуха,
загрузка POI из OpenStreetMap, персистентное хранение маршрутов в базе данных.

\section{Структура работы}

Пояснительная записка состоит из введения, четырёх глав, заключения, списка
использованных источников и приложений. В главе~2 рассматривается предметная
область; в главе~3 --- математическая модель (основной теоретический раздел);
в главе~4 --- проектирование системы; в главе~5 --- программная реализация;
в главе~6 --- тестирование. Приложения содержат фрагменты исходного кода.
